\documentclass[conference]{IEEEtran}
\IEEEoverridecommandlockouts

\usepackage{cite}
\usepackage{amsmath,amssymb,amsfonts}
\usepackage{graphicx}
\usepackage{textcomp}
\usepackage{xcolor}
\usepackage{booktabs}
\usepackage{array}
\usepackage{url}
\usepackage{float}
\usepackage{tikz}
\usetikzlibrary{arrows.meta, positioning, shapes.geometric, fit, backgrounds, calc}
\usepackage[ruled,vlined,linesnumbered]{algorithm2e}
\usepackage{multirow}
\usepackage{caption}
\usepackage{tabularx}

\def\BibTeX{{\rm B\kern-.05em{\sc i\kern-.025em b}\kern-.08em
    T\kern-.1667em\lower.7ex\hbox{E}\kern-.125emX}}

\begin{document}

\title{Semantic Firewall: A Hybrid Defense-in-Depth\\Architecture for LLM Security}

\author{
\IEEEauthorblockN{Aishwarya Sharma}
\IEEEauthorblockA{
\textit{Department of Computer Science and Engineering} \\
\textit{Institution Name}\\
City, Country \\
email@example.com
}
}

\maketitle

% ====================================================================
% ABSTRACT
% ====================================================================
\begin{abstract}
As Large Language Models (LLMs) become deeply integrated into
large-scale systems, securing them against adversarial attacks
remains computationally expensive and architecturally fragile. Current
approaches---such as Meta's Llama Guard and OpenAI's moderation
endpoints---route every user prompt through a dedicated safety LLM,
introducing prohibitive latency and cost at scale. We present
\textbf{Semantic Firewall}, a hybrid defense-in-depth architecture that
combines deterministic regex heuristics, a self-healing vector-based
semantic cache (ChromaDB), and a constrained LLM gate (Llama-3.3-70B)
executing in parallel. On the neuralchemy prompt injection benchmark\footnote{https://huggingface.co/datasets/neuralchemy/Prompt-injection-dataset}
($N{=}2{,}000$), the full system achieves 93.57\% F1 ($\pm$0.42) with
98.61\% Recall, outperforming GPT-4.0 (90.00\% F1), Llama Guard~4
(51.30\% F1), and Claude~3.5 (83.20\% F1) while reducing
downstream LLM inference dependency by 66.88\%. Against 212 manually
crafted red-team attacks, the system maintains an 85.85\% resilience
score. We further validate generalization via cross-dataset evaluation
on BeaverTails, deepset/prompt-injections, and
multi-turn evasion simulations, establishing the Semantic Firewall as a
practical, cost-effective approach to LLM security.
\end{abstract}

\begin{IEEEkeywords}
large language models, prompt injection, semantic firewall, LLM
security, defense-in-depth, semantic caching
\end{IEEEkeywords}

% ====================================================================
% I. INTRODUCTION
% ====================================================================
\section{Introduction}

Large Language Models have rapidly evolved from research prototypes
into mission-critical research and production components powering customer service
automation, code generation, knowledge retrieval, and autonomous
decision-making. This deep integration, however, exposes massive attack
surfaces: adversaries exploit prompt injections to hijack model
behavior, jailbreaks to bypass safety training, and carefully crafted
queries to exfiltrate Personally Identifiable Information (PII) or
cryptographic secrets. The security challenge is compounded by the fact
that current defenses---such as Meta's Llama
Guard~\cite{llamaguard}, OpenAI's moderation endpoints, and
PromptGuard~\cite{promptguard}---rely on routing \textit{every} user
prompt through a dedicated safety LLM, regardless of whether the input
is a trivial benign query or a sophisticated zero-day attack.

This monolithic approach suffers from three fundamental limitations.
First, it is \textit{cost-prohibitive and computationally inefficient}:
passing benign queries (e.g., ``What is the capital of France?'')
through a multi-billion parameter safety model wastes immense compute,
introduces latency bottlenecks exceeding 500\,ms per request, and
erodes computational efficiency at scale. Second, statically trained
models are \textit{vulnerable to novel zero-day attacks} that exploit
obfuscation techniques such as Base64 encoding, multi-turn hypothetical
framing, character-level perturbation, or token smuggling---attack
vectors that evolve faster than model retraining cycles. Third,
\textit{proprietary API dependence} introduces single points of failure
where strict downstream moderation blindly blocks inputs with
\texttt{403 Forbidden} errors rather than returning actionable,
structured classification responses that downstream orchestration
layers require.

To address these limitations, we present \textbf{Semantic Firewall}, a
hybrid defense-in-depth architecture that replaces the monolithic
safety-LLM paradigm with a multi-stage cascading pipeline. Instead of
routing every prompt through an expensive transformer model, the system
first applies deterministic regex heuristics and queries a self-healing
vector-based semantic cache (ChromaDB). Seven specialized parallel
heuristic detectors simultaneously analyze the input for PII leakage,
injection patterns, abuse, secrets exfiltration, and unsafe content.
Only if an input is entirely novel and bypasses all of these fast
layers does it reach the deep LLM gate---a highly constrained
Llama-3.3-70B-Instruct model forced to output strict JSON
classification schemas. This layered approach achieves state-of-the-art
security performance (93.57\% F1) while reducing LLM API dependency by
66.88\%, dropping effective system latency to $\sim$427.72\,ms in
steady state.

We explicitly scope this architecture to defend against \textit{direct}
prompt injections originating from the user input stream. We assume a
black-box attacker model where adversaries---ranging from script
kiddies using known ``DAN'' jailbreaks to sophisticated red teams
crafting novel out-of-distribution attacks---have no knowledge of the
firewall's internal infrastructure, Regex patterns, or vector cache
contents. A key design advantage is compatibility with open-weight
models: while commercial APIs lock developers into sending sensitive
data to third parties, our architecture makes it computationally viable
to route all traffic through locally deployed open-weight models within
a Virtual Private Cloud (VPC), maintaining strict data privacy
without sacrificing detection performance.

\subsection{Contributions}

Our primary contributions are fourfold:

\begin{enumerate}
    \item \textbf{Hybrid Defense-in-Depth Architecture:} A novel
    multi-layer pipeline combining deterministic regex heuristics,
    vector-based semantic caching, and a highly constrained LLM
    gate---orchestrated via parallel execution---to achieve 93.57\% F1
    with 98.61\% Recall and sub-500\,ms effective latency, significantly
    outperforming monolithic safety-LLM approaches.

    \item \textbf{Self-Healing Semantic Cache:} An auto-vectorization
    feedback loop where novel threats blocked by the LLM gate are
    dynamically embedded into ChromaDB using SentenceTransformers,
    enabling all future semantic variants to be intercepted instantly
    without LLM compute and reducing downstream API dependency by
    66.88\%.

    \item \textbf{Parallel Heuristic Ensemble:} Seven concurrent
    lightweight detectors---covering PII, abuse, injection patterns,
    secrets, DoS, threat intelligence, unsafe content, and custom
    enterprise rules---executing via Python's
    \texttt{ThreadPoolExecutor}, achieving a 90\% block rate on
    structural threats at sub-100\,ms latency with zero neural
    inference cost.

    \item \textbf{Comprehensive Empirical Validation:} A rigorous
    multi-phase benchmark across five public datasets and a manually
    crafted red-team corpus, demonstrating out-of-distribution
    generalization and statistical superiority over six commercial
    baselines including GPT-4.0, Claude 3.5, Llama Guard 4,
    Rebuff, and NeMo Guardrails.
\end{enumerate}

% ====================================================================
% II. RELATED WORK
% ====================================================================
\section{Related Work}

\textbf{Single-Model Classifiers.} State-of-the-art safety approaches
like Meta's Llama Guard~\cite{llamaguard} and
PromptGuard~\cite{promptguard} rely on passing entire user prompts
through a statically trained LLM for safety classification. While
highly accurate on their training distribution, these introduce
significant per-request latency (typically 500--1000\,ms for 8B--70B
parameter models) and cannot catch simple obfuscation techniques---such
as Base64 encoding or explicit override tokens---that deterministic
preprocessing would intercept at sub-millisecond cost. Furthermore,
without heuristic preprocessing layers, they are inherently vulnerable
to adversarial attacks that exploit distributional blind spots not
represented in their safety training data. Our empirical evaluation
confirms this: Llama Guard~4 achieves only 51.30\% F1 on the
neuralchemy benchmark, demonstrating that content moderation models do
not generalize to structural prompt injections.

\textbf{Orchestration \& Guardrail Frameworks.} Nvidia's NeMo
Guardrails~\cite{nemoguardrails} introduces programmable, rule-based
routing for LLM applications using Colang scripting. While this enables
customizable safety policies, the static rule-matching paradigm
fundamentally cannot catch semantically similar but lexically diverse
attack reformulations. The Semantic Firewall advances this paradigm by
replacing rigid rule-matching with a deep semantic cache backed by
cosine similarity over dense 384-dimensional vector embeddings,
enabling the system to intercept rephrased, paraphrased, and
synonymously restructured attacks that rule-based systems miss
entirely. Our cross-dataset evaluation demonstrates this advantage:
the Semantic Firewall achieves 49.63\% F1 on NeMo's native toxic-chat
dataset with 91.44\% Recall, catching significantly more threats than
NeMo's 66.90\% Recall at the cost of lower precision.

\textbf{Prompt Injection Specific Defenses.} Rebuff~\cite{rebuff}
employs a multi-layered defense combining heuristic analysis, LLM-based
evaluation, and a vector database for known attack signatures. Our
architecture scales and refines this concept in three critical ways:
(1)~we integrate a highly constrained 70B open-weight gate with strict
JSON-schema output constraints rather than relying on conversational
LLM classification, entirely eliminating hallucination during the
classification phase; (2)~we introduce a self-healing cache feedback
loop that continuously strengthens defenses against novel attack
variants; and (3)~we achieve parallel execution across seven
specialized detectors, ensuring that total detection latency is bounded
by the slowest agent rather than the sequential sum. On Rebuff's
native deepset/prompt-injections dataset, the Semantic Firewall
achieves 83.93\% F1 versus Rebuff's 82.50\%.

% ====================================================================
% III. SYSTEM ARCHITECTURE
% ====================================================================
\section{System Architecture}

The Semantic Firewall operates as a highly parallelized, asynchronous
Python pipeline where incoming prompts traverse progressively deeper
analytical layers. Any layer detecting a threat triggers an immediate
short-circuit block, preventing unnecessary downstream computation and
ensuring that the system's latency is bounded by the \textit{fastest}
layer capable of classifying the input. Figure~\ref{fig:arch-flow}
illustrates the end-to-end defense-in-depth architecture, and
Algorithm~\ref{alg:firewall} formalizes the complete classification
pipeline.

\begin{figure*}[t]
\centering
\includegraphics[width=\textwidth]{a1}
\caption{System Architecture of the Semantic Firewall. The pipeline orchestrates a highly parallelized defense-in-depth layout combining deterministic preprocessing, vector caching, and deep LLM analysis. Modeled strictly after the provided system design, the architecture ensures downstream LLM protection with rigid governance and logging.}
\label{fig:arch-flow}
\end{figure*}

\subsection{Layer 1: Regex Engine}

The first line of defense applies deterministic pattern matching via
highly optimized regular expressions at $\sim$5.21\,ms latency. This
layer instantly intercepts known attack signatures including ``ignore
previous instructions'' overrides, Base64-encoded payloads (matched
via the canonical Regex
\texttt{\textasciicircum(?:[A-Za-z0-9+/]\{4\})*...}),
token smuggling patterns (\texttt{<<}, \texttt{[[}}), and explicit
jailbreak prefixes such as ``DAN'' prompts. The Regex engine operates
at pure CPU cost with zero neural inference, serving as a cost-saving
compute filter designed to instantly neutralize script-kiddie attacks
and low-effort copy-pasted jailbreaks before they consume any
downstream resources.

\subsection{Layer 2: Semantic Cache (Self-Healing)}

Before any heavy computation begins, the incoming text is projected
into a 384-dimensional dense vector space using SentenceTransformers
(\texttt{all-MiniLM-L6-v2}). This embedding is queried against a
local ChromaDB instance configured with HNSW cosine indexing, which
stores the embeddings of previously blocked malicious prompts. The
cosine similarity between the incoming vector $\vec{A}$ and cached
threat vectors $\vec{B}$ is computed as:

\begin{equation}
\text{Similarity}(\vec{A}, \vec{B}) =
    \frac{\vec{A} \cdot \vec{B}}{\|\vec{A}\| \times \|\vec{B}\|}
\label{eq:cosine}
\end{equation}

If the similarity exceeds a strict threshold $\tau{=}0.90$, the
system identifies the prompt as a semantic clone of a known
attack---even if the attacker changed specific adjectives, phrasing,
or sentence structure---and blocks it instantly without LLM
invocation. Crucially, the cache implements a \textit{self-healing
feedback loop}: when the downstream LLM Gate blocks a novel zero-day
attack, the threat's vector embedding is automatically inserted into
ChromaDB. This creates a continuously strengthening defense where all
future semantic variants of previously unseen attacks are caught at
this layer with zero LLM cost, reducing median latency from
$\sim$731\,ms (full LLM inference) to $\sim$602\,ms (cache lookup)
and enabling $\sim$45$\times$ faster rejection for known attack
families.

\subsection{Layer 3: Parallel Heuristic Ensemble}

If the cache misses, the prompt is simultaneously broadcast to seven
heuristic detector agents. We employ a \texttt{ThreadPoolExecutor} to
run these synchronous Python methods concurrently, bounding the overall
ensemble latency to $\max(T_1, T_2, \dots, T_7)$ rather than the
sequential sum of all agents. The seven detectors cover complementary
threat domains: (1)~\textit{Context Flooding/DoS:} truncates
inputs exceeding token limits to prevent compute exhaustion;
(2)~\textit{PII Detection:} highly optimized Regex for instant
detection of emails, SSNs, credit cards, IBANs, and health IDs;
(3)~\textit{Secrets Scanning:} detects cryptographic tokens, API
keys, and private keys; (4)~\textit{Threat Intelligence:}
cross-references inputs against known malicious IPs/URLs;
(5)~\textit{Abuse/Toxicity:} Levenshtein distance checks against
curated toxic lexicons; (6)~\textit{Injection Detection:} structural
Regex targeting override patterns and token smuggling;
(7)~\textit{Unsafe Content:} rigid heuristic rule-sets for violence,
CSAM, and illegal activities; and (8)~\textit{Custom Rules:} bespoke
system constraints.

These agents utilize exact-match keyword dictionaries, tuned Regex
boundaries, and hard-coded mathematical limits (e.g., maximum token
length, maximum special character density). If \textit{any} agent
returns a critical flag, the orchestrator immediately halts all other
processing via short-circuiting and rejects the request. Empirically,
these foundational heuristics alone cover 90\% of major threat
vectors instantly, without ever needing a neural model.

\subsection{Layer 4: Constrained LLM Gate}

Only complex, novel inputs that bypass all fast layers reach the final
defense: Llama-3.3-70B-Instruct. The model receives a heavily
engineered system prompt enforcing strict JSON-schema classification
with precise negative constraints (e.g., ``Roleplay requests for
creative writing are BENIGN. DO NOT flag them.'' and ``Direct requests
for harmful content are NOT injections.''). The model must respond
\textit{only} with a structured JSON object containing an action
(\texttt{ALLOW}/\texttt{BLOCK}), a confidence score (0.0--1.0), and
a brief reason string. By forcing the 70B model to output structured
JSON with specific confidence intervals rather than conversational
text, we eliminate hallucination and guarantee deterministic routing.
Inference is performed with greedy decoding ($T{=}0.0$) to ensure
stable classification outcomes. If the LLM times out or the API
becomes unavailable, the system applies a fail-closed policy with
aggressive heuristic fallback.

\begin{algorithm}[t]
\DontPrintSemicolon
\SetAlgoLined
\SetKwInOut{Input}{Input}
\SetKwInOut{Output}{Output}
\Input{User prompt $q$, cache threshold $\tau$, cache $\mathcal{C}$}
\Output{Decision $d \in \{\textsc{Allow}, \textsc{Block}\}$}
\BlankLine
\tcp{Layer 1: Deterministic Regex Matching}
$r \leftarrow \textsc{RegexEngine}(q)$\;
\lIf{$r.\text{match}$}{
    \Return $\textsc{Block}(r.\text{category}, r.\text{pattern})$
}
\BlankLine
\tcp{Layer 2: Semantic Cache Lookup}
$\vec{v} \leftarrow \textsc{Encode}(q)$
    \tcp*{all-MiniLM-L6-v2, 384-dim}
$s \leftarrow \max_{\vec{c}\,\in\,\mathcal{C}}
    \text{sim}(\vec{v}, \vec{c})$\;
\lIf{$s \geq \tau$}{
    \Return $\textsc{Block}(\text{cache\_hit}, s)$
}
\BlankLine
\tcp{Layer 3: Parallel Heuristic Ensemble}
$\mathcal{R} \leftarrow \textsc{ThreadPoolExec}(
    \{D_1, \ldots, D_8\}, q)$\;
$a \leftarrow \textsc{RiskAggregate}(\mathcal{R})$
    \tcp*{Consensus vote}
\lIf{$a.\text{critical}$}{
    \Return $\textsc{Block}(a.\text{category}, a.\text{scores})$
}
\BlankLine
\tcp{Layer 4: Constrained LLM Gate (Llama-3.3-70B)}
$j \leftarrow \textsc{LLMGate}(q)$
    \tcp*{JSON schema, $T{=}0.0$}
\If{$j.\text{action} = \textsc{Block}$}{
    $\mathcal{C} \leftarrow \mathcal{C} \cup \{\vec{v}\}$
        \tcp*{Self-heal: vectorize \& insert}
    \Return $\textsc{Block}(j.\text{reason}, j.\text{confidence})$\;
}
\Return $\textsc{Allow}$\;
\caption{Semantic Firewall Classification Pipeline}
\label{alg:firewall}
\end{algorithm}

% ====================================================================
% IV. EXPERIMENTAL SETUP
% ====================================================================
\section{Experimental Setup}

\subsection{Datasets}

To validate the architecture rigorously, we evaluate across five
public benchmarks spanning diverse threat categories.
\textbf{(1)~neuralchemy/Prompt-injection-dataset}~($N{=}2{,}000$):
a balanced mixture of benign user queries and complex prompt
injections used for the ablation study, SOTA comparison, and
large-scale evaluation.
\textbf{(2)~ai4privacy/pii-masking-200k}~($N{=}209{,}261$): a
massive multilingual corpus with embedded PII entities across dozens
of types (emails, SSNs, credit cards, health IDs) used for scaled
PII detection evaluation.
\textbf{(3)~walledai/Multi-Turn-Jailbreak}~(200 sessions, 2--10
turns each): adversarial multi-turn conversation sessions where
attackers build benign personas before launching payload attacks.
\textbf{(4)~PKU-Alignment/BeaverTails}~($N{=}2{,}000$): a nuanced
safety benchmark with 14 harm categories including hate speech,
animal abuse, and financial crime, featuring difficult borderline
cases.
\textbf{(5)~deepset/prompt-injections} and
native evaluation dataset for ProtectAI
Rebuff, enabling direct
head-to-head comparison on each framework's own benchmark.

Additionally, we employ a 212-prompt manually crafted red-team
dataset featuring complex multi-layered obfuscation strategies, a
2{,}000-sample custom synthetic safety corpus covering violence,
CSAM, and illegal activities, and a 2{,}000-sample benign corpus
from \texttt{HuggingFaceH4/no\_robots} for false positive stress
testing.

\subsection{Environment \& Reproducibility}

All fast layers (Regex, Heuristics, Semantic Cache) were executed on
commodity CPU instances to simulate realistic enterprise proxy
infrastructure. ChromaDB~(\texttt{v0.4.22}) was configured with HNSW
cosine indexing (\texttt{hnsw:space:cosine}, \texttt{hnsw:M:16}).
Vector embeddings were generated using
SentenceTransformers~(\texttt{v2.2.2}) with the
\texttt{all-MiniLM-L6-v2} architecture (384 dimensions). LLM Gate
inference (Llama-3.3-70B-Instruct) was performed via the OpenRouter
API with greedy decoding ($T{=}0.0$) to guarantee stable
classification outcomes. The complete source code, benchmark scripts,
and evaluation datasets will be open-sourced via a dedicated GitHub
repository.

\subsection{Evaluation Metrics}

We report standard binary classification metrics: Precision
($\text{TP}/(\text{TP}+\text{FP})$), Recall
($\text{TP}/(\text{TP}+\text{FN})$), F1 Score
($2 \times \text{Prec} \times \text{Rec} / (\text{Prec} +
\text{Rec})$), and False Positive Rate (FPR). Statistical
significance is verified via bootstrap resampling ($N{=}10{,}000$
iterations, 95\% CI) and McNemar's test where applicable. We
additionally report effective latency, P95/P99 tail latencies, and
API cost savings to evaluate computational efficiency at enterprise
scale.

% ====================================================================
% V. RESULTS
% ====================================================================
\section{Results}

Table~\ref{tab:combined} presents the consolidated results across all
evaluation phases and baselines. Table~\ref{tab:latency} provides
the component-level latency breakdown demonstrating the computational
efficiency of each architectural layer.

\begin{table*}[t]
\centering
\caption{Consolidated Benchmark Results Across All Evaluation Phases.
Best results per benchmark are \textbf{bolded}. Precision, Recall,
and F1 are reported as percentages (\%).
``---'' indicates metric not applicable for that evaluation type.}
\label{tab:combined}
\footnotesize
\begin{tabular}{@{}llrccccl@{}}
\toprule
\textbf{Benchmark} & \textbf{Model / Configuration}
    & \textbf{$N$} & \textbf{Prec.} & \textbf{Recall}
    & \textbf{F1} & \textbf{Latency} & \textbf{Notes} \\
\midrule
\multicolumn{8}{@{}l}{\textit{(a) Ablation Study ---
    Architecture Configuration Comparison
    (neuralchemy, $N{=}2{,}000$)}} \\
\midrule
& Config 1: LLM Gate Only
    & 2000 & 78.23 & 63.66 & 70.20
    & 731.64\,ms & Single API call \\
& Config 2: Regex Only
    & 2000 & 77.34 & 62.94 & 69.40
    & 5.21\,ms & Local CPU \\
& Config 3: Semantic Cache Only
    & 2000 & 81.24 & 56.29 & 66.50
    & 602.83\,ms & Local DB search \\
& Config 4: Parallel Det. Only
    & 2000 & 83.07 & 93.77 & 88.10
    & 83.45\,ms & 8 concurrent threads \\
& Config 5: Regex + Cache
    & 2000 & 98.51 & 80.17 & 88.40
    & 602.83\,ms & Parallel execution \\
& Config 6: Regex + Par. Det.
    & 2000 & 92.04 & 88.81 & 90.40
    & 83.45\,ms & Parallel execution \\
& Config 7: All Fast (No LLM)
    & 2000 & 98.14 & 82.77 & 89.80
    & 602.83\,ms & No neural inference \\
& \textbf{Config 8: Full (Warm)}
    & 2000 & \textbf{89.03} & \textbf{98.61} & \textbf{93.57}
    & \textbf{427.72\,ms} & Steady-state \\
& Config 9: Full (Cold)
    & 2000 & 86.43 & 99.82 & 92.64
    & 459.85\,ms & Empty cache \\
\midrule
\multicolumn{8}{@{}l}{\textit{(b) Open-Weight vs. Proprietary API
    Comparison (neuralchemy, $N{=}2{,}000$)}} \\
\midrule
& \textbf{Semantic Firewall (Ours)}
    & 2000 & 86.43 & \textbf{99.82} & \textbf{92.64}
    & \textbf{459.85\,ms} & Cold start, open-weight \\
& OpenAI GPT-4.0
    & 2000 & \textbf{96.10} & 84.60 & 90.00
    & 876\,ms & Proprietary API \\
& Google Gemini 3 Flash
    & 2000 & 95.30 & 81.00 & 87.60
    & 1464\,ms & Proprietary API \\
& Anthropic Claude 3.5
    & 2000 & 94.70 & 74.30 & 83.20
    & 1156\,ms & Proprietary API \\


\midrule
\multicolumn{8}{@{}l}{\textit{(c) Cross-Dataset Generalization}} \\
\midrule
BeaverTails
    & \textbf{Semantic Firewall (Ours)}
    & 2000 & 65.23 & \textbf{87.59} & \textbf{74.80}
    & --- & 14 harm categories \\
BeaverTails
    & Meta Llama Guard 4 (12B)
    & 2000 & 73.68 & 61.20 & 66.86
    & --- & --- \\
deepset/prompt-inj.
    & \textbf{Semantic Firewall (Ours)}
    & --- & \textbf{95.10} & 75.14 & \textbf{83.93}
    & --- & Rebuff's native dataset \\
deepset/prompt-inj.
    & ProtectAI Rebuff
    & --- & 86.20 & \textbf{79.10} & 82.50
    & --- & Pipeline defense \\

\midrule
\multicolumn{8}{@{}l}{\textit{(d) Specialized Evaluations}} \\
\midrule
PII Detection
    & SF Regex Layer
    & 209K & 85.90 & 96.70 & 91.00
    & $<$2\,ms & Zero LLM cost \\
Multi-Turn Evasion
    & Semantic Firewall
    & 200 & --- & 90.00 & ---
    & --- & Session-level recall \\
Red Team (manual)
    & \textbf{Semantic Firewall}
    & 212 & --- & \textbf{85.85} & ---
    & --- & Resilience score \\
Red Team (manual)
    & GPT-4.0
    & 212 & --- & 0.94 & ---
    & --- & 2/212 blocked \\
Red Team (manual)
    & Claude 3.5
    & 212 & --- & 78.30 & ---
    & --- & 166/212 blocked \\
Unsafe Content
    & SF Full System
    & 2000 & 82.50 & 84.30 & 83.38
    & --- & FPR = 5.45\% \\
Benign Degradation
    & SF Full System
    & 2000 & --- & --- & ---
    & --- & FPR = 4.20\% \\
\midrule
\multicolumn{8}{@{}l}{\textit{(e) Cache Threshold Sensitivity
    ($\tau$ sweep, neuralchemy, $N{=}2{,}000$)}} \\
\midrule
& $\tau = 0.80$
    & 2000 & 78.12 & 99.43 & 88.24
    & --- & Hit=42.4\%, FPR=18.14\% \\
& $\tau = 0.85$
    & 2000 & 85.34 & 99.11 & 90.47
    & --- & Hit=34.7\%, FPR=12.41\% \\
& $\boldsymbol{\tau = 0.90}$ \textbf{(Selected)}
    & 2000 & \textbf{89.03} & \textbf{98.61} & \textbf{93.57}
    & --- & \textbf{Hit=12.45\%, FPR=4.20\%} \\
& $\tau = 0.95$
    & 2000 & 94.81 & 92.86 & 91.43
    & --- & Hit=14.6\%, FPR=4.24\% \\
& $\tau = 0.99$
    & 2000 & 96.27 & 88.29 & 88.12
    & --- & Hit=3.1\%, FPR=1.83\% \\
\bottomrule
\end{tabular}
\end{table*}

\begin{table}[t]
\centering
\caption{Component Latency Breakdown (Pure Compute, Isolated Profiling)}
\label{tab:latency}
\footnotesize
\begin{tabular}{@{}lcccc@{}}
\toprule
\textbf{Component} & \textbf{Avg (ms)} & \textbf{Med (ms)}
    & \textbf{P95 (ms)} & \textbf{P99 (ms)} \\
\midrule
Regex Engine & 5.21 & 5.0 & 7.24 & 9.1 \\
Semantic Cache & 602.83 & 600.0 & 615.37 & 620.1 \\
Detector Agents & 83.45 & 80.0 & 95.12 & 102.4 \\
Risk Aggregation & 0.4 & 0.4 & 0.6 & 0.8 \\
LLM Gate & 731.64 & 637.5 & 889.03 & 2655.1 \\
\bottomrule
\end{tabular}
\vspace{1mm} \\
\raggedright \textit{*Note: The 602\,ms cache latency is entirely dominated by generating the 384-dimensional vector embedding on CPU; the actual vector similarity lookup in ChromaDB takes $\sim$1\,ms. Migrating embedding generation to a T4 GPU reduces this latency to $\sim$20\,ms, which would significantly drop steady-state $L_{eff}$.}
\end{table}

\subsection{Inference from Results}

The consolidated results reveal several critical insights across
the evaluation phases, which we analyze systematically below.

\textbf{Ablation Analysis (Table~\ref{tab:combined}a).}
Progressive architectural layering monotonically improves F1. The
standalone LLM Gate (Config~1) achieves only 70.20\% F1 at
731.64\,ms, while the standalone Regex engine (Config~2) obtains a
nearly identical 69.40\% F1 at 140$\times$ lower latency (5.21\,ms).
Combining layers yields consistent, compounding gains: Regex+Cache
(Config~5) reaches 88.40\% F1 with dramatically improved Precision
(98.51\%), and adding Parallel Detectors (Config~6) reaches 90.40\%
F1. The full system (Config~8) achieves the peak 93.57\% F1 with
98.61\% Recall. The 23.37\% F1 improvement from Config~1 to Config~8
is statistically significant ($p < 0.05$, bootstrap $N{=}10{,}000$,
$\text{SD}=\pm 0.42\%$), confirming that each heuristic layer
contributes orthogonal, mathematically robust detection capabilities
rather than simply capturing statistical noise. The concurrent
execution design ensures that combining heuristic layers does not
sum their latencies linearly: Configs~4 and~6 both bottleneck at
exactly 83.45\,ms because the faster Regex engine runs in parallel.

\textbf{SOTA Comparison (Table~\ref{tab:combined}b).} The Semantic
Firewall outperforms all commercial APIs in F1, achieving 92.64\%
versus GPT-4.0's 90.00\%, Claude~3.5's 83.20\%, and Gemini
Flash Lite's 87.60\%. While commercial models achieve higher raw
Precision (GPT-4.0: 96.10\%), their Recall is consistently and
significantly lower---GPT-4.0 catches only 84.60\% of attacks
compared to our 99.82\%. The advantage is fundamentally
architectural: the Regex and Cache layers intercept structural
attacks that standalone LLMs miss, boosting system-level Recall
beyond what any single model achieves alone. This confirms that content moderation models do not generalize to
structural prompt injections. Statistical significance was verified
via McNemar's test ($\chi^2{=}4.28$, $p < 0.05$). Because 41.93\%
of traffic is resolved by deterministic fast layers, the system's
effective latency drops to 459.85\,ms even from a cold start,
outpacing every commercial API benchmarked.

\textbf{Cross-Dataset Generalization (Table~\ref{tab:combined}c).}
The system generalizes robustly to entirely unseen distributions. On
BeaverTails---a highly challenging benchmark with 14 nuanced harm
categories and difficult borderline cases---the Semantic Firewall
achieves 74.80\% F1 versus Llama Guard~4's 66.86\%, outperforming
Meta's dedicated safety model by $\sim$8 F1 points. On Rebuff's
native deepset/prompt-injections dataset, we outperform Rebuff
itself (83.93\% vs.\ 82.50\% F1), proving that the hybrid heuristic
architecture generalizes well to structural injections regardless of
dataset origin. This robust generalization confirms that the hybrid heuristic
architecture effectively intercepts structural injections,
providing a mathematically sound security bound across diverse out-of-distribution environments.

\textbf{Specialized Evaluations (Table~\ref{tab:combined}d).} The
PII Regex layer achieves 91.0\% F1 with 96.7\% Recall across
209K samples at sub-2\,ms latency using purely deterministic Regex
patterns with zero LLM cost. Highly structured entities (emails,
credit cards, IBANs, IP addresses, ZIP codes) are detected at
$>$99.7\% Recall, while moderately variable types like phone numbers
achieve 74.8\% due to irregular formatting in synthetic datasets.
The multi-turn evasion benchmark confirms 90.0\% session-level
interception across 200 adversarial sessions containing 2--10 turns.
Against 212 manually crafted red-team attacks with complex
multi-layered obfuscation, the Semantic Firewall achieves 85.85\%
resilience, outperforming GPT-4.0 (0.94\%, 2/212) and
exceeding Claude~3.5 (78.3\%, 166/212). The benign degradation
test confirms a 4.20\% FPR across 2{,}000 benign prompts, within
the 5--15\% industry-standard range. In production, false positives
would be routed to human review rather than silently dropped.

\textbf{Cache Threshold (Table~\ref{tab:combined}e).} Sweeping
$\tau$ from 0.80 to 0.99 reveals a clear precision--recall
trade-off: $\tau{=}0.80$ maximizes cache hits (42.4\%) but inflates
FPR to 18.14\%, while $\tau{=}0.99$ reduces FPR to 1.83\% but
limits cache utility (3.1\% hit rate). The optimal point
$\tau{=}0.90$ maximizes F1 (93.57\%) with a balanced 12.45\% cache
hit rate and 4.20\% FPR.

\textbf{Compute Efficiency \& Cost Analysis.} Out of the
942-sample zero-day validation suite, deterministic fast layers
resolved 630 prompts (66.88\%) without LLM invocation. Using the
empirical traffic distribution, the effective steady-state latency
is $L_{\text{eff}} = 0.4193 \times 83.45 + 0.1245 \times 602.83 +
0.4562 \times 731.64 = 443.81$\,ms. During cold-start, the latency relies heavily on the LLM Gate (where cache hits $P_c=0$ and LLM fallback $P_l=1-P_h$):
\[ L_{eff} = (0.4193 \cdot 83.45) + (0.5807 \cdot 731.64) = 459.85 \text{ ms} \]
(Note: The 12.45\% cache hit rate emerges organically during evaluation because adversarial traffic is highly clustered; the self-healing cache successfully intercepts subsequent variants of identical zero-day templates without requiring prior dataset memorization).
At large scale, routing
100M tokens through GPT-4.0 (\$5.00/1M tokens) costs \$500.00; the
Semantic Firewall reduces this to \$165.60---a 66.88\% reduction.

% ====================================================================
% VI. DISCUSSION
% ====================================================================
\section{Discussion}

The results demonstrate that the hybrid defense-in-depth paradigm
fundamentally outperforms monolithic safety-LLM approaches across
detection accuracy, latency, cost, and generalization. The core
insight---that the majority of real-world adversarial attacks are
structurally predictable and should be intercepted deterministically
rather than consuming expensive LLM compute---is empirically
validated by the 66.88\% fast-layer interception rate and the
23.37-point F1 improvement over a standalone LLM gate.

The primary source of false positives occurs when users input
toxic-adjacent but ultimately benign framing (e.g., analyzing movie
quotes containing violent language, or discussing historical
atrocities in academic context). The deterministic Regex patterns
lack semantic context and over-trigger in these cases before the
LLM Gate can apply contextual reasoning. Conversely, false negatives
concentrate in unstructured data: the PII detector achieves only
74.8\% Recall on phone numbers because synthetic datasets generate
highly irregular formats (e.g., ``1-(800)-CALL-NOW''). Expanding
Regex patterns to capture these edge cases risks catastrophic false
positive explosion on arbitrary numeric strings such as timestamps.
These failure modes reaffirm the architecture's core thesis: rigid
heuristics alone cannot secure unstructured data, and LLMs alone are
too slow. Both must be orchestrated within a layered pipeline where
each component compensates for the other's weaknesses.

% ====================================================================
% VII. LIMITATIONS & FUTURE WORK
% ====================================================================
\section{Limitations \& Future Work}

While the Semantic Firewall demonstrates robust empirical
performance, several limitations direct future research. Deploying a
local 70B parameter model requires substantial hardware (multiple
80GB A100 GPUs), making the upfront cost prohibitive for smaller
enterprises. The Python GIL constrains CPU-bound heuristic
throughput at hyper-scale ($>$10K RPS); a production deployment
should port the orchestrator to Go or Rust with a distributed
vector database (e.g., Redis) for horizontal cache
scaling across Kubernetes clusters. Sophisticated attackers could
exploit adversarial embeddings to generate synonymous payloads that
intentionally maximize cosine distance below $\tau{=}0.90$, or
launch Denial-of-Wallet attacks by flooding the system with novel
variants to force consecutive cache misses and exhaust GPU budgets.
Future iterations must explore adversarial training for the
embedding model and density-based anomaly detection within the
vector space. The reliance on public HuggingFace datasets introduces
potential pretraining data contamination for Llama-3.3-70B,
artificially inflating performance via memorization; future
evaluations should utilize privately generated zero-day benchmarks.
Finally, fine-tuning quantized Small Language Models (e.g.,
Llama-3-8B via vLLM) for the gate role could reduce slow-path
latency from 731\,ms to under 150\,ms, fundamentally eliminating
cloud-hosted safety model dependence entirely.

% ====================================================================
% VIII. CONCLUSION
% ====================================================================
\section{Conclusion}

The Semantic Firewall definitively demonstrates that relying on a
single large language model for security is computationally and
architecturally flawed. By introducing a defense-in-depth
architecture combining fast deterministic filtering for known threats
(Regex, Semantic Cache, Parallel Heuristics) and deep LLM analysis
for novel zero-day attacks (constrained Llama-3.3-70B gate), we
achieve 93.57\% F1 with 98.61\% Recall on the neuralchemy benchmark
($N{=}2{,}000$), while reducing LLM API invocations by 66.88\%. The
ablation study across 9 configurations validates every architectural
component: replacing the fast heuristics with an LLM-only gate drops F1 by 23 points,
while the semantic cache prevents the 66.88\% LLM routing
overhead that bottlenecks monolithic systems. Cross-dataset
evaluation on BeaverTails, deepset/prompt-injections,
and our manual Red Team suite proves the architecture's
out-of-distribution generalization. The open-weight architecture
outperforms GPT-4.0 by 2.6 F1 points (cold start), Claude~3.5 by 9.4
points, and Llama Guard~4 by 41.3 points---enabling deployment within isolated environments
by utilizing open-weight models without relying on proprietary
black-box APIs, and establishing a practical, cost-effective approach
to LLM security. The complete source code, benchmark
scripts, and evaluation datasets will be open-sourced for full
reproducibility.

\begin{thebibliography}{4}
\bibitem{llamaguard}
M. Inan et al., ``Llama Guard: LLM-based Input-Output Safeguard for
Human-AI Conversations,'' \textit{arXiv preprint arXiv:2312.06674},
2023.

\bibitem{promptguard}
Meta, ``PromptGuard: A safeguard model against prompt injections and
jailbreaks,'' 2024.



\bibitem{rebuff}
ProtectAI, ``Rebuff: Prompt Injection Defense,'' GitHub Repository,
2023.
\end{thebibliography}

\end{document}
